\SourceLeaf{206}Man kann also den Vortrag so einkleiden: Die Gleichung
\[
\left.\begin{array}{c} S \\ \Sigma \end{array}\right\} = M\pm h+i+u
\]
lehrt, dass gleiche Stellungen der Blase, bei unge\"andertem Werth von $i$, gleiche Werthe von $u$ anzeigen. Nun sind die Werthe von $u$ in I und II entgegengesetzt; hat [die Blase] also in II dieselbe Stellung wie in I, so ist $u'=u''=0$.

Im entgegengesetzten Fall schliesst man, dass $u$ nicht $=0$ war und jetzt den entgegengesetzten Werth vom vorigen hatte. Bringt man also durch stetiges Schrauben an 7 das $u$ wieder auf den vorigen Werth, so geht es dabei durch 0, und das halbe Schrauben zur\"uck bringt 0 hervor. Dann ist also die Linie 1\,2 oder 2\,1 oder 3\,4 horizontal, aber die Blase hat sich ihrerseits wieder gleichfalls der 2$^{\mathrm{ten}}$ Stellung um die H\"alfte gen\"ahert. Man \"andert nun $i$ durch Schrauben an 3, bis $\left.\substack{S\\\Sigma}\right\} = M\pm h$ wird.

\begin{center}
\textsc{Bemerkungen.}
\end{center}

Die erste der vorstehenden Notizen findet sich auf dem hintern Schmutzblatte des Werkes: \emph{Mathematische Abhandlungen} von Jacob Wilhelm Heinrich \textsc{Lehmann}, Zerbst 1829 und stammt vermuthlich aus diesem Jahre. Die Notiz [2] steht auf einem einzelnen Blatte; die Zeit ihrer Abfassung l\"asst sich nicht genau angeben, doch darf man vermuthen, dass sie zwischen 1820 und 1830 f\"allt. Die Notiz [3] hat \textsc{Gauss} in einem Handbuche verzeichnet, h\"ochst wahrscheinlich im M\"arz 1832, denn sie bezieht sich auf \textsc{Gauss}’ Brief an \textsc{Bolyai} vom 6. M\"arz 1832. Die Notiz [4] beruht auf einer m\"undlichen Mittheilung von \textsc{Gauss} an \textsc{Lübsen}, der im Jahre 1830 bei ihm Vorlesungen h\"orte (vergl. \textsc{Lübsen}, \emph{Ausf\"uhrliches Lehrbuch der Analysis}, Hamburg 1853, S. 171). Endlich ist die Notiz [5] auf einem Zettel verzeichnet, der nicht datirt ist, jedoch deutet die Handschrift auf eine sp\"atere Periode aus \textsc{Gauss}’ Leben, etwa die Zeit zwischen 1840--1850; wahrscheinlich hat \textsc{Gauss} die Notiz bei der Vorbereitung f\"ur eine Vorlesung \"uber praktische Astronomie niedergeschrieben. F\"ur ihre Unterbringung an dieser Stelle ist lediglich ihr Inhalt massgebend gewesen.

Dass \textsc{Gauss} sich schon sehr fr\"uh mit der Definition der Ebene besch\"aftigt hat, zeigt die bereits S. 162 dieses Bandes abgedruckte Stelle aus seinem Tagebuche vom 28. Juli 1797.

Die beiden Figuren in der Notiz [1] sind dem Texte hinzugef\"ugt worden.

\hfill \textsc{Stäckel.}

